\documentclass[11pt,a4paper]{article}

\usepackage[T1]{fontenc}
\usepackage[utf8]{inputenc}
\usepackage[margin=1in]{geometry}
\usepackage{lmodern}
\usepackage{microtype}
\usepackage{xcolor}
\usepackage{graphicx}
\usepackage{booktabs}
\usepackage{tabularx}
\usepackage{array}
\usepackage{enumitem}
\usepackage{caption}
\usepackage{subcaption}
\usepackage{float}
\usepackage{parskip}
\usepackage{hyperref}
\usepackage{titlesec}
\usepackage{fancyhdr}
\usepackage{tcolorbox}

\hypersetup{
  colorlinks=true,
  linkcolor=black,
  urlcolor=blue!60!black,
  pdftitle={Interactive Wall Calendar Component Report},
  pdfauthor={Aryan Bokolia}
}

\definecolor{reportnavy}{HTML}{1F2937}
\definecolor{reportteal}{HTML}{1F766B}
\definecolor{reportsoft}{HTML}{EFF8F6}
\definecolor{reportline}{HTML}{D7DDE5}
\definecolor{reportmuted}{HTML}{64748B}

\pagestyle{fancy}
\fancyhf{}
\fancyhead[L]{Interactive Wall Calendar Component}
\fancyhead[R]{Official Submission Report}
\fancyfoot[C]{\thepage}
\setlength{\headheight}{14pt}

\titleformat{\section}
  {\Large\bfseries\color{reportnavy}}
  {\thesection}{0.6em}{}

\titleformat{\subsection}
  {\normalsize\bfseries\color{reportnavy}}
  {\thesubsection}{0.6em}{}

\newcommand{\detailrow}[2]{\textbf{#1} & #2 \\}

\newcommand{\reportimage}[2][0.98\linewidth]{%
  \IfFileExists{#2}{%
    \includegraphics[width=#1]{#2}%
  }{%
    \fbox{%
      \parbox[c][5.4cm][c]{0.92\linewidth}{%
        \centering
        \textbf{Image not uploaded yet} \\[0.5em]
        Upload this file to Overleaf: \texttt{#2}
      }%
    }%
  }%
}

\begin{document}

\pagenumbering{gobble}

% Overleaf note:
% Upload this file along with the following image assets for the final version:
% public/report/hero-section.png
% public/report/month-section.png
% public/report/notes-section.png
% public/report/reminders-section.png

\begin{titlepage}
  \centering

  {\Large\textsc{National Institute of Technology Karnataka, Surathkal}\par}
  \vspace{0.5cm}
  {\large Department of Computer Science and Engineering\par}
  \vspace{1.4cm}

  {\Huge\bfseries Interactive Wall Calendar Component\par}
  \vspace{0.4cm}
  {\Large Enhanced Frontend Engineering Challenge Submission\par}
  \vspace{1.2cm}

  \begin{tcolorbox}[
    colback=reportsoft,
    colframe=reportline,
    boxrule=0.6pt,
    arc=4mm,
    width=0.92\textwidth
  ]
    \centering
    \textbf{Submission Summary} \\[0.5em]
    This report documents an enhanced version of the original frontend engineering challenge.
    The final project extends the brief into a production-minded interactive calendar experience
    with polished UX, responsive design, accessibility improvements, reminders, notes, and
    refined interface behavior.
  \end{tcolorbox}

  \vfill

  \renewcommand{\arraystretch}{1.35}
  \begin{tabularx}{0.9\textwidth}{>{\raggedright\arraybackslash}p{0.28\textwidth}X}
    \detailrow{Submitted By}{Aryan Bokolia}
    \detailrow{Student ID}{241CS111}
    \detailrow{Program}{B.Tech. in Computer Science and Engineering}
    \detailrow{Institution}{National Institute of Technology Karnataka, Surathkal}
    \detailrow{Location}{Karnataka, India}
    \detailrow{Submission Type}{Official Frontend Engineering Challenge Report}
    \detailrow{Project Stack}{Next.js, TypeScript, Tailwind CSS, Framer Motion, date-fns}
    \detailrow{Submission Date}{\today}
  \end{tabularx}

  \vfill
\end{titlepage}

\clearpage
\pagenumbering{arabic}

\section{Project Overview}

This project began as a frontend engineering challenge to build an interactive wall calendar
component. The implementation was deliberately extended beyond the base requirements to
demonstrate stronger frontend engineering judgment across interaction design, accessibility,
responsiveness, visual polish, and real-world usability.

The final product combines a wall-calendar-inspired hero layout, an interactive month grid,
a persistent notes workflow, reminder scheduling, motion design, theme support, and robust
edge-case handling. The goal was not only to complete the assignment, but to produce a
submission that feels product-ready and professionally presented.

\section{Submission Context and Objectives}

\begin{itemize}[leftmargin=1.4em]
  \item Build a visually distinctive wall-calendar interface rather than a generic calendar widget.
  \item Implement date range selection with intuitive interaction patterns, including drag selection.
  \item Add practical product features such as notes, reminders, notifications, and quick actions.
  \item Ensure responsive behavior across desktop and mobile form factors.
  \item Improve accessibility through keyboard navigation, visible focus states, and readable contrast.
  \item Present the project as an official, recruiter-ready engineering submission.
\end{itemize}

\section{Interface Showcase}

\begin{figure}[H]
  \centering
  \begin{subfigure}[t]{0.48\textwidth}
    \centering
    \reportimage{public/report/hero-section.png}
    \caption{Hero section with wall-calendar styling and editorial framing.}
  \end{subfigure}
  \hfill
  \begin{subfigure}[t]{0.48\textwidth}
    \centering
    \reportimage{public/report/month-section.png}
    \caption{Month view with selected range, quick actions, and date markers.}
  \end{subfigure}
  \caption{Primary calendar experience.}
\end{figure}

\begin{figure}[H]
  \centering
  \reportimage[0.95\linewidth]{public/report/notes-section.png}
  \caption{Notes workflow tied to the currently selected date range.}
\end{figure}

\begin{figure}[H]
  \centering
  \reportimage[0.95\linewidth]{public/report/reminders-section.png}
  \caption{Reminders workflow with timezone-aware targeting and notification status.}
\end{figure}

\section{Implemented Features}

\subsection{Core Features}

\begin{itemize}[leftmargin=1.4em]
  \item Wall calendar layout with a hero image and structured month-view interface.
  \item Current month display using a seven-column calendar grid.
  \item Highlighting for the current day and selected range.
  \item Fully responsive layout for desktop, tablet, and mobile devices.
\end{itemize}

\subsection{Interaction and Productivity Features}

\begin{itemize}[leftmargin=1.4em]
  \item Click-based date range selection.
  \item Drag-to-select interaction inspired by booking and scheduling platforms.
  \item Quick actions including Today, Clear Selection, Select This Week, and month navigation.
  \item Notes system with create, edit, delete, and localStorage persistence.
  \item Reminders with browser notifications and in-app fallback alerts.
  \item Timezone-aware reminder labeling and scheduling behavior.
\end{itemize}

\subsection{Visual and UX Enhancements}

\begin{itemize}[leftmargin=1.4em]
  \item Framer Motion animations for load states, date selection, and month transitions.
  \item Light and dark mode support with persisted user preference.
  \item Dynamic accent styling influenced by the hero image.
  \item Important-date indicators for today, notes, reminders, and special dates.
  \item Toast-style feedback for note and reminder actions.
\end{itemize}

\subsection{Accessibility and Edge-Case Handling}

\begin{itemize}[leftmargin=1.4em]
  \item Keyboard navigation using arrow keys, Enter, and Space.
  \item Visible focus states and improved screen-reader clarity.
  \item Validation for invalid selections, empty note content, and incomplete reminder input.
  \item Graceful behavior for reversed ranges, same-day selection, and denied notification permissions.
\end{itemize}

\section{Design and Engineering Decisions}

\subsection{Why the wall calendar concept was chosen}

The wall calendar format gives the project a stronger identity than a standard date picker.
It creates a more memorable submission and makes the interface feel like a designed product
rather than a default utility component.

\subsection{Why localStorage was used}

The challenge was frontend-focused, so localStorage was selected to keep the application
self-contained, easy to run, and simple for reviewers to evaluate without needing backend
infrastructure, authentication, or cloud services.

\subsection{How UX was approached}

The interface was refined with an emphasis on smooth interaction flow, reduced friction, clear
feedback, and realistic usage patterns. This includes responsive spacing, drag selection, sticky
or full-width workflow sections, action feedback, mobile-friendly controls, and restrained motion.

\section{Technology Stack}

\begin{tcolorbox}[
  colback=reportsoft,
  colframe=reportline,
  boxrule=0.6pt,
  arc=3mm
]
\begin{itemize}[leftmargin=1.4em]
  \item Next.js 16 with App Router
  \item TypeScript
  \item Tailwind CSS
  \item Framer Motion
  \item date-fns
  \item localStorage for client-side persistence
\end{itemize}
\end{tcolorbox}

\section{Future Improvements}

\begin{itemize}[leftmargin=1.4em]
  \item Backend integration for cloud persistence and cross-device synchronization.
  \item User authentication and personalized calendar state.
  \item Cloud-based reminders with more reliable background delivery.
  \item Shared planning, collaboration, and multi-user calendar workflows.
\end{itemize}

\section{Conclusion}

This submission represents an enhanced and more production-minded interpretation of the original
frontend engineering challenge. The project demonstrates not only visual implementation skill,
but also thoughtful UX decisions, accessibility improvements, maintainable architecture, and
practical product-oriented feature design.

Overall, the work reflects a strong frontend engineering focus with attention to usability,
interaction quality, and professional presentation.

\end{document}
